\documentclass[12pt]{article}
\usepackage{amsmath, amssymb, amsfonts}
\usepackage{amsthm}
\usepackage{geometry}
\usepackage{booktabs}
\geometry{margin=1in}

\newtheorem{proposition}{Proposition}

\title{CES Labor Income and the Effect of Part-Time Training on Saving}
\author{}
\date{}

\begin{document}
\maketitle

\noindent This note generalizes the two-period saving problem in which labor income is multiplicative, $wzx$, to a CES aggregator $wf(x,z)$. The object of interest is how the sign of $\rho$ changes the result that part-time training crowds out saving at low income and crowds in saving at high income.

\section{Setup}

The household solves
\begin{equation}
    \max_{a'}\ \ln(c)+\beta\mathbb{E}_{z'}[\ln(c')]
\end{equation}
subject to
\begin{equation}
    c+a'=(1+r)a+n\,w f(x,z),
    \qquad
    c'=(1+r')a'+w'f(x',z'),
\end{equation}
where
\begin{equation}
    f(x,z)
    =
    \bigl[\alpha_1 x^{1-\rho}+\alpha_2 z^{1-\rho}\bigr]^{1/(1-\rho)},
    \qquad
    \alpha_1,\alpha_2>0.
    \label{eq:ces-f}
\end{equation}
The elasticity of substitution between skill and the productivity shock is $\sigma=1/\rho$. Special cases of~\eqref{eq:ces-f}:
\begin{itemize}
    \item $\rho=0$: perfect substitutes, $f=\alpha_1 x+\alpha_2 z$;
    \item $\rho\to 1$: Cobb--Douglas;
    \item $\rho>1$: gross complements ($\sigma<1$);
    \item $\rho<0$: $f_z$ decreases in skill (super-substitutes). Here $1-\rho>1$, so $f$ is convex in $(x,z)$ and $\sigma=1/\rho<0$ is no longer an elasticity of substitution; this case lies outside the standard concave CES range and is included only as a formal benchmark.
\end{itemize}
The baseline technology $xz$ is \emph{not} nested in~\eqref{eq:ces-f}: $xz$ is homogeneous of degree~2, whereas~\eqref{eq:ces-f} is constant returns to scale.

Compare part-time training $(n=1,e=e_L)$ with no training $(n=1,e=0)$. Training scales skill from $x$ to $tx$ with $t>1$. Use the same normalization as in the multiplicative analysis: $(1+r)=(1+r')=w=w'=1$, period-1 shock $z=1$, and $\ln z'\sim\mathcal{N}(0,\sigma_z^2)$. Write
\begin{equation}
    \mu\equiv\mathbb{E}[z']=e^{\sigma_z^2/2},
    \qquad
    \Sigma\equiv\operatorname{Var}(z')=e^{\sigma_z^2}(e^{\sigma_z^2}-1).
\end{equation}
The budgets become
\begin{equation}
    c+a'=a+f(x,1),
    \qquad
    c'=a'+f(tx,z').
\end{equation}
Let $y_1\equiv a+f(x,1)$. The Euler equation is
\begin{equation}
    \frac{1}{y_1-a'}
    =
    \beta\,\mathbb{E}_{z'}\!\left[\frac{1}{a'+f(tx,z')}\right].
\end{equation}

\section{CES derivatives and income-risk exposure}

Write $A=\alpha_1 x^{1-\rho}+\alpha_2 z^{1-\rho}$, so $f=A^{1/(1-\rho)}$. Then
\begin{equation}
    f_z=\alpha_2\Bigl(\frac{f}{z}\Bigr)^{\rho},
    \qquad
    f_x=\alpha_1\Bigl(\frac{f}{x}\Bigr)^{\rho},
\end{equation}
and therefore
\begin{equation}
    \frac{\partial f_z}{\partial x}
    =
    \alpha_1\alpha_2\rho\,\frac{f^{2\rho-1}}{x^{\rho}z^{\rho}}.
    \label{eq:fz-x}
\end{equation}
Equation~\eqref{eq:fz-x} is the key comparative static:
\begin{equation}
    \operatorname{sign}\!\left(\frac{\partial f_z}{\partial x}\right)
    =
    \operatorname{sign}(\rho).
\end{equation}
Part-time training raises the sensitivity of labor income to $z'$ if and only if $\rho>0$. If $\rho<0$, additional skill \emph{insures} against the shock. If $\rho=0$, that sensitivity is constant.

\section{Second-order approximate saving}

Match the multiplicative derivation: evaluate expected period-2 labor income at $z'=\mu$ and its variance by the delta method.\footnote{\label{fn:approx}This treats $\mathbb{E}[f(tx,z')]\approx f(tx,\mu)$, which is exact when $f$ is linear in $z'$ (the multiplicative case) but under CES drops the mean correction $\tfrac12 f_{zz}(tx,\mu)\Sigma$, which is of the same order in $\Sigma$ as the retained precautionary term. A consistent second-order expansion of the Euler equation adds
$\Gamma(t,x)\equiv -f_{zz}(tx,\mu)\Sigma/[2(1+\beta)]$ to~\eqref{eq:astar-ces}, where $f_{zz}=-\alpha_1\alpha_2\rho\,f^{\rho}x^{1-\rho}z^{-\rho-1}/A$, so $\operatorname{sign}(f_{zz})=-\operatorname{sign}(\rho)$. The term vanishes for $\rho=0$. For $\rho>1$ its $t$-derivative is $O(x^{\rho})=o(\varphi_t)$ as $x\to0$, leaving low-income crowd-out unchanged, but as $x\to\infty$ it is of the same order as $\varphi_t$, with $\partial_t f_{zz}/\varphi_t\to\rho(\rho-1)/\mu^{2}$; Part~3 of Proposition~\ref{prop:dadt-ces} accounts for this. For $\rho<0$ the sign of $\partial_t f_{zz}$ is ambiguous at intermediate $x$ but the term is $O(x^{\rho})=o(1)$ as $x\to\infty$, so Part~1 holds asymptotically. Parts~1 and~2 below use the simpler form~\eqref{eq:astar-ces} to parallel the multiplicative analysis.} Define
\begin{equation}
    \varphi(t,x)\equiv f(tx,\mu),
    \qquad
    \psi(t,x)\equiv f_z(tx,\mu)=\alpha_2\bigl(\varphi/\mu\bigr)^{\rho}.
\end{equation}
Then $\mathbb{E}[f(tx,z')]\approx\varphi$ and $\operatorname{Var}(f(tx,z'))\approx\psi^2\Sigma$. Linearizing the Euler equation around the certainty-equivalent point gives
\begin{equation}
    a'^{\star}(x,a;t)
    =
    \underbrace{\frac{\beta y_1-\varphi}{1+\beta}}_{\text{CE}}
    +
    \underbrace{\frac{\kappa\,\varphi^{2\rho}}{\beta(y_1+\varphi)}}_{\text{Precautionary}},
    \qquad
    \kappa\equiv \alpha_2^2\Sigma\,\mu^{-2\rho}.
    \label{eq:astar-ces}
\end{equation}
The multiplicative formulas are the special case $\varphi=tx\mu$ and $\kappa\varphi^{2\rho}=t^2x^2\Sigma$.

Differentiate~\eqref{eq:astar-ces} with respect to $t$. Since $\varphi_t=\alpha_1(\varphi/(tx))^{\rho}\,x>0$,
\begin{equation}
    \frac{\partial a'^{\star}}{\partial t}
    =
    \varphi_t
    \left[
        -\frac{1}{1+\beta}
        +\frac{\kappa}{\beta}H'(\varphi)
    \right],
    \qquad
    H(\varphi;y_1)\equiv\frac{\varphi^{2\rho}}{y_1+\varphi},
    \label{eq:dadt-ces}
\end{equation}
where $H'\equiv\partial H/\partial\varphi$ and
\begin{equation}
    H'(\varphi;y_1)
    =
    \frac{\varphi^{2\rho-1}\bigl[2\rho y_1+(2\rho-1)\varphi\bigr]}{(y_1+\varphi)^2}.
    \label{eq:Hprime}
\end{equation}
The first term in the bracket of~\eqref{eq:dadt-ces} is the certainty-equivalent crowd-out: training raises expected period-2 income. The second term is the precautionary response. Its sign is $\operatorname{sign}(H')$ and is therefore governed by $\rho$. Define, as in the multiplicative analysis,
\begin{equation}
    \Delta(x,a;t)
    \equiv
    a'^{\star}(x,a;t)-a'^{\star}(x,a;1)
    =
    \int_1^t \frac{\partial a'^{\star}}{\partial u}(x,a;u)\,du.
\end{equation}

\section{How the sign of $\rho$ changes the result}

\begin{proposition}\label{prop:dadt-ces}
    Under the approximation~\eqref{eq:astar-ces}, the effect of part-time training on saving is governed by $\rho$ as follows.
    \begin{enumerate}
        \item If $\rho\le 0$, then $\partial a'^{\star}/\partial t<0$ and $\Delta(x,a;t)<0$ for all $x>0$, $a\ge 0$, and $t>1$. Training crowds out saving at every income level. The high-income crowd-in disappears.
        \item If $0<\rho<1$, then
        \begin{equation}
            \frac{\Delta(x,a;t)}{x}
            \;\xrightarrow{x\to\infty}\;
            -\frac{\alpha_1^{1/(1-\rho)}(t-1)}{1+\beta}<0.
            \label{eq:rho-lt-1-limit}
        \end{equation}
        Consumption smoothing dominates for sufficiently large $x$, so high-income crowd-in does not survive.
        \item Let $\rho>1$, $a>0$, and use the consistent second-order expansion of footnote~\ref{fn:approx}. Write
        \begin{equation}
            \varphi_{\infty}=\alpha_2^{1/(1-\rho)}\mu,
            \qquad
            y_{1\infty}=a+\alpha_2^{1/(1-\rho)}.
        \end{equation}
        If
        \begin{equation}
            \frac{\alpha_2^2\,H'(\varphi_{\infty};y_{1\infty})}{\beta\,\mu^{2\rho}}
            -\frac{\rho(\rho-1)}{2(1+\beta)\mu^{2}}
            >
            \frac{1}{(1+\beta)\,\Sigma},
            \label{eq:ces-sigma-star}
        \end{equation}
        then $\Delta<0$ for small $x$ and $\Delta>0$ for large $x$. The qualitative multiplicative result is restored.
    \end{enumerate}
\end{proposition}

\begin{proof}
    \emph{Part 1.} If $\rho<0$, then $2\rho y_1+(2\rho-1)\varphi<0$, so $H'<0$ by~\eqref{eq:Hprime}. Both terms in the bracket of~\eqref{eq:dadt-ces} are negative. Hence $\partial a'^{\star}/\partial t<0$ and $\Delta<0$. Skill substitutes for luck: training raises expected period-2 income and reduces exposure to $z'$.

    If $\rho=0$, then $f=\alpha_1 x+\alpha_2 z$, $\varphi=\alpha_1 tx+\alpha_2\mu$, and $\psi=\alpha_2$ (variance independent of $t$). Differentiating~\eqref{eq:astar-ces} gives
    \[
        \frac{\partial a'^{\star}}{\partial t}
        =
        -\frac{\alpha_1 x}{1+\beta}
        -\frac{\alpha_2^2\Sigma\,\alpha_1 x}{\beta(y_1+\varphi)^2}<0.
    \]
    Again $\Delta<0$ for all $x$.

    \emph{Part 2.} If $0<\rho<1$, skill dominates as $x\to\infty$:
    \[
        \varphi(t,x)\sim\alpha_1^{1/(1-\rho)}\,tx,
        \qquad
        y_1\sim\alpha_1^{1/(1-\rho)}\,x.
    \]
    The certainty-equivalent gap is of order $x$. The precautionary gap is of order $x^{2\rho-1}$. Because $2\rho-1<1$, the ratio $\Delta/x$ converges to the certainty-equivalent limit~\eqref{eq:rho-lt-1-limit}. If $\Sigma$ is very large, $\Delta$ can be positive at intermediate $x$, but it turns negative again as $x\to\infty$.

    \emph{Part 3.} If $\rho>0$, then $f_z$ increases in skill, so training raises period-2 income risk. For $\rho\ge 1/2$, $H'>0$ always and the two channels oppose each other.

    Now take $\rho>1$ and $a>0$, and add the mean correction $\Gamma(t,x)=-f_{zz}(tx,\mu)\Sigma/[2(1+\beta)]$ of footnote~\ref{fn:approx} to~\eqref{eq:astar-ces}, so that
    \[
        \frac{\partial a'^{\star}}{\partial t}
        =
        \varphi_t\left[-\frac{1}{1+\beta}+\frac{\kappa}{\beta}H'(\varphi;y_1)\right]
        -\frac{\Sigma}{2(1+\beta)}\,\partial_t f_{zz}(tx,\mu).
    \]
    As $x\to 0$, $\varphi\sim\alpha_1^{1/(1-\rho)}\,tx\to 0$, $y_1\to a>0$, and $H'(\varphi;y_1)\to 0$ because $2\rho-1>0$. Moreover $\alpha_1(tx)^{1-\rho}/A\to 1$, so $f_{zz}(tx,\mu)\sim-\alpha_2\rho\,\varphi^{\rho}\mu^{-\rho-1}$ and $\partial_t f_{zz}=O(x^{\rho})=o(\varphi_t)$ since $\varphi_t=O(x)$ and $\rho>1$. Thus
    \[
        \frac{\partial a'^{\star}}{\partial t}\approx -\frac{\varphi_t}{1+\beta}<0,
    \]
    so $\Delta<0$ at low $x$. As $x\to\infty$, the shock is the bottleneck:
    \[
        \varphi\to\varphi_{\infty}=\alpha_2^{1/(1-\rho)}\mu,
        \qquad
        y_1\to y_{1\infty}=a+\alpha_2^{1/(1-\rho)},
        \qquad
        \varphi_t\sim\alpha_1\varphi_{\infty}^{\rho}\,t^{-\rho}x^{1-\rho}\to 0.
    \]
    Then $H'(\varphi_{\infty};y_{1\infty})>0$ because $\rho>1/2$. For the mean correction, $A\to\alpha_2\mu^{1-\rho}$ gives $f_{zz}(tx,\mu)\sim-\alpha_1\rho\,\mu^{-2}\varphi_{\infty}^{\rho}(tx)^{1-\rho}$, hence
    \[
        \partial_t f_{zz}(tx,\mu)\sim\alpha_1\rho(\rho-1)\,\mu^{-2}\varphi_{\infty}^{\rho}\,t^{-\rho}x^{1-\rho}
        =\frac{\rho(\rho-1)}{\mu^{2}}\,\varphi_t .
    \]
    Substituting,
    \[
        \frac{\partial a'^{\star}}{\partial t}
        \sim
        \varphi_t\left[-\frac{1}{1+\beta}\Bigl(1+\frac{\rho(\rho-1)\Sigma}{2\mu^{2}}\Bigr)+\frac{\kappa}{\beta}H'(\varphi_{\infty};y_{1\infty})\right].
    \]
    With $\kappa=\alpha_2^2\Sigma\mu^{-2\rho}$, the bracket is positive if and only if~\eqref{eq:ces-sigma-star} holds. Because $\varphi(u,x)\to\varphi_{\infty}$ for every $u\in[1,t]$, the bracket has the same sign along the whole path of integration in $\Delta$, so $\Delta>0$ for large $x$. Since $\varphi_t\to 0$, $\Delta\to 0$ from above.
\end{proof}

\paragraph{Interpretation of condition~\eqref{eq:ces-sigma-star}.}
Unlike the multiplicative case, a large $\Sigma$ alone does not deliver high-income crowd-in. The left-hand side of~\eqref{eq:ces-sigma-star} must be positive: the precautionary gain from higher exposure to $z'$, $\alpha_2^2 H'(\varphi_\infty;y_{1\infty})/(\beta\mu^{2\rho})$, must exceed $\rho(\rho-1)/[2(1+\beta)\mu^{2}]$. The second term arises because $f$ is concave in $z'$ when $\rho>0$, so expected period-2 income lies $\tfrac12 f_{zz}\Sigma$ below $\varphi$. Training shrinks this concavity penalty, raising expected income and adding a crowd-out channel that also scales with $\Sigma$; the two $\Sigma$-terms therefore compete regardless of how large $\Sigma$ is.

\paragraph{Cobb--Douglas ($\rho=1$).}
If $\alpha_1+\alpha_2=1$, then $f=x^{\alpha_1}z^{\alpha_2}$. Increments in the mean and in the variance are both of order $x^{\alpha_1}$. The sign of $\Delta$ is a horse race between $\Sigma$ and $\alpha_1$, as in the multiplicative model (which is closer to $\alpha_1$ near one).

\section{Discussion}

The multiplicative technology $txz'$ makes $\operatorname{Var}(c')$ of order $x^2$. The certainty-equivalent and precautionary gaps therefore remain the same order as $x\to\infty$, and a single crossing can occur when $\Sigma/\mu$ is large enough. That knife-edge scaling is why the original proposition can deliver high-income crowd-in.

Under the CRS CES~\eqref{eq:ces-f}, the same qualitative conclusion requires $\rho>1$ (complements) together with condition~\eqref{eq:ces-sigma-star}, which is not implied by a large $\Sigma$ alone. Even then the high-income crowd-in is weaker than in the multiplicative model: because $\varphi_t=O(x^{1-\rho})\to 0$, $\Delta$ is positive but vanishes as $x\to\infty$, whereas under $txz'$ the gap $\Delta$ grows linearly in $x$. If $\rho\le 0$, both channels crowd out saving and the sign flip disappears.

\begin{center}
\begin{tabular}{lll}
\toprule
$\rho$ & Exposure $f_z$ vs.\ skill & Effect of part-time training on saving \\
\midrule
$\rho<0$ & falls & $\Delta<0$ for all $x$ (no crowd-in) \\
$\rho=0$ & flat & $\Delta<0$ for all $x$ (no crowd-in) \\
$0<\rho<1$ & rises, slowly & CE wins as $x\to\infty$; high-$x$ crowd-in fails \\
$\rho=1$ & rises with income & horse race in $\Sigma$; original logic possible \\
$\rho>1$ & rises, strongly & original result: low-$x$ out, high-$x$ in, if~\eqref{eq:ces-sigma-star} holds \\
\bottomrule
\end{tabular}
\end{center}

\end{document}
